        \btabularx
        \hline
        \multicolumn{4}{|l|}{\bf queue mode attributes} \\
        \hline
        depth          & 3PL     & int     & queue buffer depth                          \\         
        mindepth       & 3PL     & int     & minimum queue buffer depth                  \\         
        buffersize     & 3PL     & int     & queue buffer depth                          \\         
        minbuffersize  & 3PL     & int     & minimum queue buffer depth                  \\         
        readdomain     & 3PL     & clock   & force the read (output) clock domain        \\    
        writedomain    & 3PL     & clock   & force the write (input) clock domain        \\    
        domain         & 3PL     & clock   & force both the read and write clock domains \\    
        alu            & 3PL     & boolean & use an ALU if appropriate                   \\       
        dsp            & 3PL     & boolean & use a DSP block                             \\       
        fifo           & 3PL     & boolean & use a hardware FIFO                         \\       
        queuereg       & 3PL     & boolean & extra output register to improve timing     \\       
        continuous     & 3PL     & boolean & no startup dependence                       \\   
        ignoremodules  & 3PL     & boolean & all reads treated as if in same module      \\   
        readonly       & 3PL     & boolean & set to read only                            \\
        queuefieldzero & netlist & boolean & clear unassigned fields                     \\
        \hline
        \btabularxend
        
        {\bf These attributes are case-insensitive.}
        
        {\bf depth} sets the depth of the queue buffer. The default depth is 2
        for a synchronous queue and 16 for an asynchronous queue.
        The values allowed depend on the FPGA family.

        \begin{tabular}{| l | l | l |} \hline
                                & using cmemory  & using rmemory \\ \hline \hline
        Spartan-II, Spartan IIE & 16             &  256 to  4096 \\ \hline
        Spartan-3               & 16             &  512 to 16384 \\ \hline
        Virtex                  & 16 to  32      &  256 to  4096 \\ \hline
        Virtex-E                & 16 to  32      &  256 to  4096 \\ \hline
        Virtex-II               & 16 to  64      &  512 to 16384 \\ \hline
        Virtex-II               & 16 to  64      &  512 to 16384 \\ \hline
        Virtex 4                & 16             &  512 to 16384 \\ \hline
        Virtex 5                & 16 to 128      & 1024 to 32768 \\ \hline                                        
        \end{tabular}
        
        A depth of 0 may also be set. This is a special case, creating an
        unbuffered queue. The use of unbuffered queues is very restricted -
        see section \autorefph{sec:unbufferedqueues}.

        {\bf buffersize} alternative older attribute name for {\bf depth}.

        {\bf mindepth} sets the depth of the queue buffer but rounds the
        supplied depth up to the next allowed value.

        {\bf minbuffersize} alternative older attribute name for {\bf mindepth}.

        {\bf readdomain} forces the read (output) clock domain to the specified
        clock domain. {\bf writedomain} forces the write (input) clock domain to
        the specified clock domain.
        {\bf domain} forces both the read (output) clock domain and the write
        (input) clock domain to the specified clock domain. 
        If {\bf null} is assigned, the current clock domain
        is selected. If the variable already has a clock domain as a result of
        having been assigned or evaluated then that clock domain is compared with
        the intended domain. If the assigned clock domain agrees with that already
        in place then there is no further action, otherwise a fatal error occurs.
        
        {\bf alu} specifies that assignments of expressions containing
        arithmetic operators '+' and '-' as the final or only operation can be
        combined into a single common adder/subtracter using CLBs. If the {\bf alu}
        attribute is false it can override the global directive {\bf ALU} to
        suppress the use of an ALU for an individual static variable.
        
        The {\bf alu} attribute can save logic and incurs no time penalty for the
        operations drawn in to the adder/subtracter, but all other assignments
        suffer an increase in propagation time since assigned values must now
        traverse the shared adder/subtracter. This tradeoff must be considered
        in each case.
        
        The {\bf alu} attribute cannot be added or changed once the variable has been
        assigned.
        
        The behaviour of the {\bf dsp} attribute depends on the FPGA -
        \blist
        \item   For Virtex-5 and Virtex-6 it has a similar effect to the
                {\bf alu} option in that it draws arithmetic and bitwise logic
                operations in assignments to the variable into a single ALU,
                but using a DSP block rather than CLBs.
        \item   For other FPGAs the attribute has no effect.
        \blend
        
        The {\bf dsp} attribute will not use a Virtex-5 or Virtex-6 DSP block as an
        ALU if the word width exceed 48 bits (multiplier operands are restricted
        to 25 by 18 bits).
        
        Note that for Virtex-5 or Virtex-6 if the {\bf alu} attribute is true but not
        the {\bf dsp} attribute then arithmetic operators '+' and '-' as the final or
        only operation will be combined into a single common adder/subtracter
        using CLBs.
        
        The {\bf fifo} attribute when true specifies that the queue buffer will
        use inbuilt block RAM FIFO logic if possible. It is ignored if the FPGA
        does not have this logic or the queue variable is asynchronous and uses
        the read or write status outputs. This attribute cannot be used if the
        queue is type "null", has the default depth of 2, or has a default depth
        which such that CLB RAM is used (depth \verb'<=' 16). In these cases an
        informative message is produced and the attribute is ignored. FIFO
        hardware can be enabled globally by setting directive {\bf global.FIFO}
        in which case this attribute can be assigned false to counteract the
        directive, i.e. the attribute when set overrides the directive. The FIFO
        directive is false by default.
        
        Use of the FIFO hardware is a compromise between speed and logic
        usage. The FIFO hardware saves CLB gates and flip-flops, which would
        otherwise be used, but has greater path delays which may cause
        failure to meet timing constraints.
        
        The {\bf queuereg} attribute when true causes an extra pipelined output
        register to be added. This reduces the clock-to-output time. The queue
        buffer logic incorporates the extra buffer with no change to
        functionality.

        {\bf Continuous} true specifies that if possible assignment will occur
        as soon as the FPGA is operational after configuration. This can
        increase operating speed by eliminating clock-enable logic. There are
        two situations where this can occur -
        \blist
        \item   the queue variable has a non-conditional assignment - the
                assignment operates continuously following configuration loading
                regardless of any 3PL startup mechanisms.
        \item   a {\bf when} statement has only unconditional assignments to
                variables in its nested code block(s) - the assignments are
                enabled directly from the {\bf when} test expression (or its
                inverse) and are not controlled from execution threads
                initiated by the 3PL startup mechanisms.
        \blend
        see {\bf start()}(\autorefph{sec:ipstart}).
        Note that if directive "continuous" is defined as type log with
        value true at the end of immediate execution, all queue variables
        will be treated as though they had the continuous attribute true.

        {\bf ignoremodules} true specifies that all queue reads will be
        treated as though from within the same module. Thus regardless of
        which module a queue read is contained within it will pop the queue
        when executed and the availability of the next datum will be seen
        within all modules in which the queue is read. This is distinct from
        the default behaviour of queue reads from multiple modules where a
        queue read from each relevant module must be executed before the next
        datum (if any) is seen as available to all the modules. This attribute
        is ignored for unbuffered queues.
        
        The purpose of this option is to allow code within multiple modules
        to read queues while looking after exclusive access where necessary
        by means of explicit user code. This can also be done also by ensuring
        that all the queue reads are contained within one module but this
        restriction can be inconvenient and may result in more obscure code.
        
        {\bf readonly} true specifies that from now on this variable cannot
        be assigned to. This attribute cannot be assigned false.
        
        {\bf queuefieldzero} false inhibits the clearing of unassigned fields or
        array members when a queue variable is assigned. The default is {\bf true}.
        If this attribute is made false for a queue variable which has a
        non-primitive type and an incomplete assignment is made, i.e. not all
        struct fields or array members are assigned, the unassigned elements
        for the datum will be undefined instead of explicitly zero. The value
        in making this attribute false is that it saves logic resources and
        reduces the delays on some data paths. Where a queue variable has a primitive
        type, or the whole variable is assigned in one statement or in multiple statements
        a par or sync block, this attribute has no effect.
